\documentclass{beamer}
% COMSC-3013 Computer Architecture shared Beamer preamble.
% Week decks live one directory below weeks/ and input this file as:
%   % COMSC-3013 Computer Architecture shared Beamer preamble.
% Week decks live one directory below weeks/ and input this file as:
%   % COMSC-3013 Computer Architecture shared Beamer preamble.
% Week decks live one directory below weeks/ and input this file as:
%   \input{../_shared/beamer-preamble.tex}
% Keep the course self-contained. Change visual/build conventions here,
% not by forking one-off week themes.

\usetheme{Boadilla}
\usecolortheme{seahorse}
\setbeamertemplate{navigation symbols}{}
\setbeamertemplate{footline}[frame number]

\usepackage{amsmath}
\usepackage{booktabs}
\usepackage{graphicx}
\usepackage{listings}
\usepackage{pgfpages}

\lstset{
  basicstyle=\ttfamily\small,
  columns=fullflexible,
  keepspaces=true,
  showstringspaces=false,
  breaklines=true
}

% Speaker notes remain hidden in the ordinary student build. Define \shownotes
% before inputting the deck to create an instructor copy; see weeks/README.md.
\ifdefined\shownotes
  \setbeameroption{show notes on second screen=right}
\fi

\newcommand{\ArchTitleSlide}[4]{%
  % #1 week label, #2 title, #3 central question, #4 Wednesday prediction
  \begin{frame}
    \titlepage
  \end{frame}
  \begin{frame}{The machine question}
    \Large #3

    \vspace{1.2em}
    \normalsize\textbf{Prediction to carry into Wednesday:} #4
  \end{frame}
}

% Keep the course self-contained. Change visual/build conventions here,
% not by forking one-off week themes.

\usetheme{Boadilla}
\usecolortheme{seahorse}
\setbeamertemplate{navigation symbols}{}
\setbeamertemplate{footline}[frame number]

\usepackage{amsmath}
\usepackage{booktabs}
\usepackage{graphicx}
\usepackage{listings}
\usepackage{pgfpages}

\lstset{
  basicstyle=\ttfamily\small,
  columns=fullflexible,
  keepspaces=true,
  showstringspaces=false,
  breaklines=true
}

% Speaker notes remain hidden in the ordinary student build. Define \shownotes
% before inputting the deck to create an instructor copy; see weeks/README.md.
\ifdefined\shownotes
  \setbeameroption{show notes on second screen=right}
\fi

\newcommand{\ArchTitleSlide}[4]{%
  % #1 week label, #2 title, #3 central question, #4 Wednesday prediction
  \begin{frame}
    \titlepage
  \end{frame}
  \begin{frame}{The machine question}
    \Large #3

    \vspace{1.2em}
    \normalsize\textbf{Prediction to carry into Wednesday:} #4
  \end{frame}
}

% Keep the course self-contained. Change visual/build conventions here,
% not by forking one-off week themes.

\usetheme{Boadilla}
\usecolortheme{seahorse}
\setbeamertemplate{navigation symbols}{}
\setbeamertemplate{footline}[frame number]

\usepackage{amsmath}
\usepackage{booktabs}
\usepackage{graphicx}
\usepackage{listings}
\usepackage{pgfpages}

\lstset{
  basicstyle=\ttfamily\small,
  columns=fullflexible,
  keepspaces=true,
  showstringspaces=false,
  breaklines=true
}

% Speaker notes remain hidden in the ordinary student build. Define \shownotes
% before inputting the deck to create an instructor copy; see weeks/README.md.
\ifdefined\shownotes
  \setbeameroption{show notes on second screen=right}
\fi

\newcommand{\ArchTitleSlide}[4]{%
  % #1 week label, #2 title, #3 central question, #4 Wednesday prediction
  \begin{frame}
    \titlepage
  \end{frame}
  \begin{frame}{The machine question}
    \Large #3

    \vspace{1.2em}
    \normalsize\textbf{Prediction to carry into Wednesday:} #4
  \end{frame}
}

\title{Week 16 Monday}
\subtitle{Farkle + Machine Learning: What Did the Compute Buy?}
\author{COMSC-3013}
\date{}
\begin{document}
\begin{frame}\titlepage\end{frame}
\ArchQuestionSlide{What does it cost a machine to make a better Farkle decision, and when is that extra cost actually worth paying?}{Choose two fixed strategies. Predict effectiveness, preparation cost, operating cost, and whether the extra computation will be worth it under one declared objective.}

\begin{frame}{The shortcut we are testing}
\begin{itemize}
\item More computation is not automatically a better architecture.
\item Faster hardware is not automatically the right machine.
\item The workload, objective, and constraint decide what matters.
\end{itemize}
\note{Frame Week 16 as payoff and judgment, not a new technical unit.}
\end{frame}

\begin{frame}{Two different axes}
\begin{itemize}
\item Algorithmic effort: threshold, learned table, bounded rollout.
\item Execution substrate: the machine and execution path underneath the fixed workload.
\item Change both at once and causal interpretation gets muddy.
\end{itemize}
\end{frame}

\begin{frame}{Four currencies}
\begin{itemize}
\item Effectiveness: what did the strategy accomplish?
\item Preparation: what did it cost before play?
\item Operation: what did it cost while playing?
\item Machine/environment: what actually executed the work?
\end{itemize}
\end{frame}

\begin{frame}{Pay now or pay later?}
\begin{itemize}
\item Learned strategy: preparation/training cost first, cheap decisions later.
\item Rollout strategy: little preparation, repeated simulation during decisions.
\item Heuristic: tiny cost may still be rational under a bounded objective.
\end{itemize}
\note{Use amortization language without opening a new math unit.}
\end{frame}

\begin{frame}{Throughput is evidence, not scripture}
\begin{itemize}
\item Wall-clock measurements have noise.
\item Repeat the same fixed workload.
\item Report median, minimum, and maximum games/second.
\end{itemize}
\end{frame}

\begin{frame}{Correctness before speed}
\begin{itemize}
\item Balanced starters and raw denominators stay visible.
\item Deterministic seeds keep the software comparison stable.
\item A faster wrong benchmark does not win.
\end{itemize}
\end{frame}

\begin{frame}{A GPU in the chassis proves nothing}
\begin{itemize}
\item Required Week 16 path is native Python CPU.
\item Physical accelerator presence is not accelerator use.
\item A future accelerator lane must prove dispatch and equivalence.
\end{itemize}
\end{frame}

\begin{frame}{Your prediction}
\begin{itemize}
\item Pick two fixed strategies.
\item Predict effectiveness.
\item Predict where each pays its computational cost.
\item Name the objective that will determine your verdict.
\end{itemize}
\end{frame}

\begin{frame}{Wednesday}
\begin{itemize}
\item Hold game rules, seed, workload, and machine lane fixed.
\item Run repeated evidence.
\item Read JSON/CSV, not just terminal output.
\item Ask: what did the extra computation buy?
\end{itemize}
\end{frame}

\begin{frame}{Scope}
\begin{itemize}
\item No Checkpoint 4.
\item No new Machine Dossier layer.
\item No required GPU, cloud, Kubernetes, or paid AI.
\item One playful workload, one architectural judgment.
\end{itemize}
\end{frame}

\begin{frame}{Final question}
\centering\Large When is ``more compute'' actually worth paying for?
\end{frame}
\end{document}
